\section{作る機能から使う場へ}
\sectionkicker{MULTIMODAL}
\begin{wideflow}
{\Large\bfseries 作る機能から使う場へ}\par
\smallskip
{\color{SurveyMuted}音・像・映像の動きは模型名の列ではなく、使える場と移行の記録として読む。}\par
\end{wideflow}

\noindent 8月20日のAdobeの記事は、Fireflyにおける音楽生成、音声生成、効果音生成の広範な提供を中心に置く。音楽は商用に安全な licensed の楽曲、音声はAdobeとElevenLabsの2系統、効果音は動きや間、勢いに合わせる内容として描かれる。「商用に安全」という言い回しは提供側の表現であり、利用条件や選択肢の同等性は実際の約款に譲る。品質や網羅の主張も提供側の説明であり、独自の試しはない。 \cite{w2026w34w34eventc091}\par\medskip

\noindent Googleの廃止頁は、Imagen 4の3つの一般提供の受け口が2026年8月17日に止まり、gemini-3.1-flash-imageに置き換わるとする。Firebase側の移行案内は、全てのImagen模型が開発者向けAPIと応用基盤を横断して同日に止まり、道具の形に破壊的な変わり目があると述べる。Vertex AI側の面には6月30日という別の日付が載るため、確かめの対象はGemini側の8月17日に置き、Vertex側の範囲は未確定として残す。当日の時刻の刻みは取得文にない。 \cite{w2026w34w34eventc018}\par\medskip

\noindent Pikaの8月14日の頁は、映像から音を起こすSoundtrack、文から効果を起こすSFX、読み上げのSpeech、楽曲全体のMusicという4模型の家族を示し、秒・分あたりの料金をAPI Club経由で載せる。競合比の倍率の主張は単位の揃わない提供側の報告であり、独自の確かめはない。8月18日の模型ごとの記事は家族の発表の敷衍であり、4件の追加発表として数えない。Alibabaの工房表は、高速な映像生成とされるwan3.0-video-primeを8月20日の段に載せ、8月6日の段と区別する。日付の刻みのみであり、能力の言い回しは提供側の表現である。後のRunway側の統合とは混同しない。DeepSeekの8月21日の変更記録は、実験的な映像理解の模型DeepSeek-V4-Flash-Vision-ExpをAPI上で出し、純文の力は公式のV4-Flash並みで、型名の指定で使えるとする。同等性や数値の主張は提供側の説明であり、手法や外部の評価は取得文にない。 \cite{w2026w34w34eventc020,w2026w34w34eventgapalibabawan30modelstudio,w2026w34w34eventc022}\par\medskip

\begin{claimboundary}[読解上の境界]
この節の読解上の境界: 'Commercially safe' is Adobe framing; license terms and ElevenLabs-option parity require the actual terms, not captured here. / Quality/coverage claims are vendor-stated without independent testing in captured bytes. / Date-level precision only. / Capability descriptors are Alibaba-stated. / Cost-multiple claims (9x/20x/10x/2x vs named rivals) are Pika-reported on non-normalized units; not independently verified. / Aug 18 per-model posts elaborate the Aug 14 family launch; they are not four additional launches. / Vertex AI surfaces carry a different documented date (June 30) for the same IDs; the Gemini-API Aug 17 date is the verified claim, Vertex scope unresolved. / Exact UTC cutoff moment on Aug 17 is unspecified in captured bytes. / Capability-parity and any benchmark figures are vendor-stated; methodology and independent evaluation are absent from captured bytes. / commercial-safety claims remain vendor framing / capability quality is not inferred from availability alone / migration/deprecation event and new-model launch are distinct identities
\end{claimboundary}
